\subsection{The Determinant of a $2\times2$ Matrix}

\begin{definitionbox}
For
\[
A=\begin{pmatrix}a&b\\c&d\end{pmatrix},
\]
the determinant is
\[
\boxed{\det A=ad-bc}.
\]
We also write
\[
\begin{vmatrix}a&b\\c&d\end{vmatrix}=ad-bc.
\]
\end{definitionbox}

The products use the two diagonals, but the order is essential: the second product is subtracted from the first.

\begin{examplebox}[A nonzero determinant]
For
\[
A=\begin{pmatrix}3&1\\2&5\end{pmatrix},
\]
we obtain
\[
\det A=3\cdot5-1\cdot2=13.
\]
\end{examplebox}

\begin{examplebox}[A zero determinant]
For
\[
B=\begin{pmatrix}1&2\\3&6\end{pmatrix},
\]
\[
\det B=1\cdot6-2\cdot3=0.
\]
The cancellation will later be explained by the proportional rows.
\end{examplebox}

\begin{interpretationbox}[Position matters]
Exchanging the rows changes
\[
\det\begin{pmatrix}1&2\\3&4\end{pmatrix}=-2
\]
into
\[
\det\begin{pmatrix}3&4\\1&2\end{pmatrix}=2.
\]
The same entries in different positions can give a different determinant.
\end{interpretationbox}
